\documentclass[11pt]{article}

\usepackage[margin=1.15in]{geometry}
\usepackage{amsmath,amssymb,amsthm}
\usepackage{booktabs}
\usepackage{microtype}
\usepackage{xcolor}
\usepackage[colorlinks=true,linkcolor=blue!50!black,citecolor=blue!50!black,urlcolor=blue!50!black]{hyperref}

\newtheorem{theorem}{Theorem}[section]
\newtheorem{lemma}[theorem]{Lemma}
\newtheorem{corollary}[theorem]{Corollary}
\newtheorem{proposition}[theorem]{Proposition}
\newtheorem{remark}[theorem]{Remark}
\newtheorem{question}[theorem]{Question}
\theoremstyle{definition}
\newtheorem{definition}[theorem]{Definition}

\newcommand{\N}{\mathbb{N}}
\newcommand{\Z}{\mathbb{Z}}
\newcommand{\Q}{\mathbb{Q}}
\newcommand{\R}{\mathbb{R}}
\newcommand{\Ztwo}{\Z_2}
\newcommand{\Qtwo}{\Q_2}
\newcommand{\code}[1]{\texttt{#1}}

\title{The Word-Complexity Ladder for Collatz Divergence:\\
  a Machine-Verified Rung Zero and a Mahler-Value\\
  Reduction of the Automatic Rung}
\author{M.~Sharpe\thanks{\texttt{msharpe@irendi.com}. The mathematics,
  the Lean formalization, and this text were produced in collaboration
  with Claude (Anthropic). The Collatz conjecture remains open, and
  nothing in this paper claims otherwise.}}
\date{June 12, 2026}

\begin{document}
\maketitle

\noindent\fbox{\parbox{\dimexpr\linewidth-2\fboxsep}{\textbf{Status note (August 2026).} The rung-one question left open here --- the irrationality of the 2-adic Mahler value $F_2(8/9)$ --- has since been settled and machine-verified: no orbit has the itinerary $\sigma^\infty(1)$ or any shift of it (\code{lean/Collatz/Rung1.lean}), and the truncation exponent $0.946$ discussed in Section~\ref{sec:gap} is a barrier for truncation only; a degree-one Pad\'e form suffices. See the companion paper \emph{Rung One of the Word-Complexity Ladder} (\code{paper/rung1.pdf}). Rung zero, the reduction, and the calibration below stand as written.}}
\medskip

\begin{abstract}
Order potential divergent orbits of the Terras map ($T(n) = n/2$ for
even $n$, $(3n+1)/2$ for odd $n$) by the complexity of their parity
itineraries. We prove the floor of this ladder formally: a divergent
orbit's itinerary is never eventually periodic (machine-verified in
Lean~4, axiom-free). We then reduce the next rung to a single
arithmetical statement. For the fixed point $w$ of the substitution
$1 \mapsto 110$, $0 \mapsto 011$ --- an aperiodic automatic word of
parity density $2/3$, inside the supercritical band where divergent
orbits must live --- the generating function $F(z) = \sum w_i z^i$
satisfies the Mahler functional equation
$F(z) = (z+z^2)/(1-z^3) + (1-z^2)\,F(z^3)$, and the unique $2$-adic
integer whose itinerary is $w$ equals $-10 + \tfrac{5}{9} F_2(8/9)$,
where $F_2$ denotes $2$-adic evaluation. Consequently, if the $2$-adic
Mahler value $F_2(8/9)$ is irrational, no integer has itinerary $w$.
The archimedean value $F(8/9) \in \R$ is transcendental by Mahler's
classical method; the $2$-adic value is the open case, and we prove
that it is genuinely beyond naive methods: the truncation
approximations to $F_2(8/9)$ have irrationality exponent exactly
$(3\log 2)/(2\log 3) = 0.9464 < 1$, so no Liouville-style argument
can decide the question --- and the deficiency constant is, once
again, the critical density of the $3x+1$ problem. The same reduction
applies uniformly to every constant-length substitution with constant
block weight in the supercritical band, where double (real and
$2$-adic) convergence of the relevant series is automatic.
\end{abstract}

\section{Introduction}

The Collatz conjecture asserts that iterating the Terras map
\[
T(n) \;=\;
\begin{cases}
n/2 & n \text{ even},\\
(3n+1)/2 & n \text{ odd},
\end{cases}
\]
from any positive integer eventually reaches $1$
\cite{lagarias-overview}. Its negation has two faces: a nontrivial
cycle, or a divergent orbit. This paper concerns the divergent face,
through the lens of symbolic dynamics. The \emph{itinerary} of $n$ is
the infinite binary word
\[
w(n) \;=\; \bigl(T^{i}(n) \bmod 2\bigr)_{i \ge 0},
\]
the diary of parity choices the orbit makes. Two facts, both formal in
the companion repository \cite{companion}, frame everything below.
First, the itinerary determines the integer completely: distinct
positive integers have distinct (in fact everywhere eventually
differing) itineraries --- the infinite-precision form of Terras'
parity bijection \cite{terras}. Second, a divergent orbit cannot have
an arbitrary itinerary: its odd-step density is pinned, in the liminf,
to at least the critical value
\[
\delta_c \;=\; \frac{\log 2}{\log 3} \;=\; 0.630929\ldots
\]
(the density floor of \cite{companion}; intuitively, fewer odd steps
than that and the halvings win).

\textbf{The ladder.} Order itineraries by word-theoretic complexity:
eventually periodic words; automatic words (fixed points of
constant-length substitutions and their codings); morphic and Sturmian
words; and beyond. Excluding each class from the divergent repertoire
is a theorem target. The conjecture --- on its divergence side --- is
the top of the ladder; each rung is finite, honest work. This paper
establishes the floor and gives the second rung a precise,
classical shape.

\textbf{Results.}
\begin{enumerate}
\item \emph{Rung zero, machine-verified}
(Section~\ref{sec:rungzero}). A divergent orbit's itinerary is not
eventually periodic. Formalized in Lean~4
(\code{divergent\_itinerary\_aperiodic}, file \code{Ladder.lean}
of \cite{companion}), depending only on Lean's standard foundations;
no axioms, no \code{sorry}.
\item \emph{The automatic rung is one Mahler value}
(Sections~\ref{sec:mahler}--\ref{sec:reduction}). For the fixed point
$w$ of $\sigma: 1 \mapsto 110,\ 0 \mapsto 011$, the unique $2$-adic
integer with itinerary $w$ is $x(w) = -10 + \tfrac59 F_2(8/9)$, where
$F(z) = \sum_i w_i z^i$ satisfies
$F(z) = (z+z^2)/(1-z^3) + (1-z^2)F(z^3)$. Hence if $F_2(8/9) \notin
\Q$, no integer realizes $w$. The word $w$ is aperiodic
(Lemma~\ref{lem:aperiodic}), $F$ is transcendental over $\Q(z)$
(Lemma~\ref{lem:fatou}), the real value $F(8/9)$ is transcendental by
Mahler's classical theorem (Theorem~\ref{thm:arch}), and the statement
generalizes uniformly to all constant-length, constant-block-weight
substitutions in the supercritical band
(Proposition~\ref{prop:class}).
\item \emph{The cheap route is provably insufficient}
(Section~\ref{sec:gap}). The truncations of the defining series give
rational approximations to $F_2(8/9)$ with irrationality exponent
exactly $(3\log 2)/(2\log 3) < 1$; rational numbers admit
approximations of exponent $1$. No argument that uses only the
convergence speed of the series can decide irrationality. The
deficiency ratio is $\delta_c$ again, in disguise, and it caps the
naive exponent below $1$ throughout the supercritical band ---
the same coincidence that makes the density theory critical makes the
transcendence question non-trivial.
\end{enumerate}

We emphasize the modesty of the target: the automatic rung is a
countable, measure-zero family of itineraries, and closing it would
leave the conjecture wide open (the companion no-go theorem
\cite{companion} shows no bounded-memory argument reaches the top of
the ladder). Its interest is structural --- it is the first rung where
the divergence problem meets a developed branch of transcendence
theory through an exact interface, with every reduction step proved
and most of them machine-checked.

\section{Itineraries, the correction term, and realizations}
\label{sec:realization}

Write $T^{i}$ for the $i$-th iterate, $j_T(n)$ for the number of odd
values among $n, T(n), \ldots, T^{T-1}(n)$ (the odd-step count), and
abbreviate the itinerary prefix $w_{<T} = (w_0, \ldots, w_{T-1})$,
$w_i = T^i(n) \bmod 2$. The dynamics admits an exact linear form
(\code{terras\_exact\_form} in \cite{companion}): for all $n, T$,
\begin{equation}\label{eq:exact}
2^T \, T^{T}(n) \;=\; 3^{j_T} n \;+\; d(w_{<T}),
\end{equation}
where the correction term $d$ depends only on the parity word, via
the closed form (\code{dcoef\_closed})
\begin{equation}\label{eq:closed}
d(w_{<T}) \;=\; \sum_{\substack{i < T \\ w_i = 1}}
  2^{i}\, 3^{\#\{i' > i \,:\, i' < T,\ w_{i'} = 1\}}.
\end{equation}
Both are theorems of the formal development; we use them as given.

Now let $w$ be an \emph{infinite} binary word. Since $3$ is invertible
in $\Ztwo$, \eqref{eq:exact} reads $n \equiv -3^{-j_T} d(w_{<T})
\pmod{2^T}$, and from \eqref{eq:closed},
\[
3^{-j_T}\, d(w_{<T}) \;=\; \sum_{\substack{i < T \\ w_i = 1}}
2^{i}\, 3^{-(1 + \nu_i)}, \qquad
\nu_i := \#\{i' < i : w_{i'} = 1\},
\]
\emph{exactly} (the exponent in \eqref{eq:closed} is
$j_T - \nu_i - 1$). The right-hand side is a partial sum of a series
whose $i$-th term has $2$-adic absolute value $2^{-i}$; it therefore
converges in $\Ztwo$.

\begin{definition}\label{def:realization}
The \emph{realization} of an infinite word $w$ is
\[
x(w) \;:=\; -\sum_{i \,:\, w_i = 1} 2^{i}\, 3^{-(1+\nu_i)}
\;\in\; \Ztwo .
\]
\end{definition}

\begin{proposition}\label{prop:realization}
A positive integer $n$ has itinerary $w$ if and only if $x(w) = n$
in $\Ztwo$. In particular at most one integer realizes any given word.
\end{proposition}

\begin{proof}
If $n$ has itinerary $w$ then $n \equiv -3^{-j_T}d(w_{<T}) \pmod{2^T}$
for every $T$ by \eqref{eq:exact}, and the right-hand sides converge
to $x(w)$; hence $n = x(w)$. Conversely if $x(w) = n$ then $n$ agrees
with $-3^{-j_T} d(w_{<T})$ modulo $2^T$ for every $T$; by the parity
bijection (the first $T$ parities of an orbit are determined by the
starting point mod $2^T$, \code{parity\_of\_modEq} in
\cite{companion}, after \cite{terras}) and an induction on $T$, the
itinerary of $n$ is $w$.
\end{proof}

Excluding a word $w$ from the divergent repertoire therefore means
proving $x(w) \notin \Z_{\ge 1}$ --- a statement about one $2$-adic
number per word.

\section{Rung zero}\label{sec:rungzero}

\begin{theorem}[machine-verified]\label{thm:rungzero}
Let $n$ have unbounded orbit: for every $B$ there is $t$ with
$T^{t}(n) > B$. Then the itinerary of $n$ is not eventually periodic.
\end{theorem}

\begin{proof}[Proof sketch]
If the itinerary satisfies $w_{s+t} = w_{s+P+t}$ for all $t \ge 0$,
then the two tail points $u = T^{s}(n)$ and $v = T^{s+P}(n)$ have
identical itineraries at every time; by the infinite-precision parity
bijection (two integers with everywhere-equal itineraries agree modulo
$2^k$ for every $k$, hence are equal --- \code{eq\_of\_itinerary\_eq})
we get $u = v$. The orbit is then eventually periodic, so it is
bounded by the maximum of its first $s + P + 1$ values, contradicting
unboundedness. The full proof is
\code{divergent\_itinerary\_aperiodic} in \code{Ladder.lean} of
\cite{companion}; it depends only on \code{propext},
\code{Classical.choice}, \code{Quot.sound}.
\end{proof}

Rung zero is elementary, but it sets the pattern: a complexity class
of words is excluded by showing its realizations are never positive
integers (here, eventually periodic words realize the rational
$2$-adic points of the dynamics --- e.g.\ $x((110)^\infty) = -5$ ---
and the positive-integer ones among them are the bounded orbits).

\section{The automatic rung: one functional equation}
\label{sec:mahler}

From here on, $w$ is the fixed point of
\[
\sigma:\quad 1 \mapsto 110, \qquad 0 \mapsto 011,
\]
starting from $w_0 = 1$:
$w = 110\,110\,011\,110\,110\,011\,011\,110\,110\cdots$. Both images
have length $3$ and contain exactly two $1$s, so $w$ has parity
density exactly $2/3 > \delta_c$: a word in the supercritical band,
admissible (as far as the density theory knows) as the itinerary of a
divergent orbit. The substitution structure is equivalent to the
\emph{kernel relations}
\begin{equation}\label{eq:kernel}
w_{3i} = w_i, \qquad w_{3i+1} = 1, \qquad w_{3i+2} = 1 - w_i
\qquad (i \ge 0),
\end{equation}
read off from $\sigma(1) = 110$, $\sigma(0) = 011$.

\begin{lemma}\label{lem:aperiodic}
$w$ is not eventually periodic.
\end{lemma}

\begin{proof}
Suppose $w_{i+P} = w_i$ for all $i \ge N$, with $P \ge 1$ minimal
admissible. If $3 \mid P$, write $P = 3Q$: for $i \ge N$,
\eqref{eq:kernel} gives $w_{i+Q}$ $= w_{3(i+Q)} = w_{3i + P} = w_{3i}
= w_i$ for all $i \ge N$, so $Q < P$ is also an eventual period;
hence we may assume $3 \nmid P$. Consider positions
$i \equiv 1 \pmod 3$, $i \ge N$; there $w_i = 1$ by
\eqref{eq:kernel}. If $P \equiv 1 \pmod 3$ then $i + P \equiv 2
\pmod 3$, so $1 = w_{i+P} = 1 - w_{(i+P-2)/3}$, i.e.\
$w_{t + (P-1)/3} = 0$ for every $t \ge (N-1)/3$ (writing $i = 3t+1$).
Then $w$ is eventually all $0$s, contradicting $w_{3i+1} = 1$. If
$P \equiv 2 \pmod 3$ then $i + P \equiv 0 \pmod 3$, so
$1 = w_{i+P} = w_{(i+P)/3}$, i.e.\ $w_{t + (P+1)/3} = 1$ for all large
$t$: $w$ is eventually all $1$s, contradicting
$w_{3i+2} = 1 - w_i = 0$ at the infinitely many $i$ with $w_i = 1$.
\end{proof}

\begin{proposition}[the functional equation]\label{prop:feq}
The generating function $F(z) = \sum_{i \ge 0} w_i z^i \in \Z[[z]]$
satisfies
\[
F(z) \;=\; \frac{z + z^2}{1 - z^3} \;+\; (1 - z^2)\, F(z^3).
\]
\end{proposition}

\begin{proof}
Group the sum in blocks of three and apply \eqref{eq:kernel}:
\[
F(z) = \sum_{i \ge 0} z^{3i}\bigl(w_i + z + (1 - w_i)z^2\bigr)
     = (z + z^2)\sum_{i\ge0} z^{3i} + (1 - z^2)\sum_{i\ge0} w_i z^{3i}.
\qedhere
\]
\end{proof}

This is the classical Mahler setup \cite{mahler, nishioka}: degree
$d = 3$, coefficients in $\Q(z)$, iteration $z \mapsto z^3$. (It is
also verified coefficient-wise to degree $900$ by the scripts of
\cite{companion}, as are all identities below; see
Section~\ref{sec:repro}.)

\begin{lemma}\label{lem:fatou}
$F$ is transcendental over $\Q(z)$.
\end{lemma}

\begin{proof}
The coefficients of $F$ lie in $\{0,1\}$. If $F$ were rational, its
coefficients would eventually satisfy a fixed linear recurrence with
constant coefficients; the sequence of state windows (the $D$ most
recent coefficients, for $D$ the recurrence depth) takes values in
the finite set $\{0,1\}^D$, so some window repeats, and forward
determinism makes the coefficient sequence eventually periodic ---
contradicting Lemma~\ref{lem:aperiodic}. So $F$ is irrational; by
Fatou's theorem \cite{fatou} (a power series with integer
coefficients and radius of convergence $1$ is either rational or
transcendental over $\Q(z)$; see also \cite[Ch.~6]{as}),
$F$ is transcendental over $\Q(z)$.
\end{proof}

\begin{theorem}[the evaluation identity]\label{thm:eval}
In $\Qtwo$ (where $|8/9|_2 = 1/8 < 1$ makes all series below
converge),
\[
x(w) \;=\; -10 \;+\; \tfrac{5}{9}\, F_2(8/9),
\]
where $F_2(8/9) := \sum_{i \ge 0} w_i (8/9)^i \in \Ztwo$ denotes the
$2$-adic evaluation.
\end{theorem}

\begin{proof}
Each block of three letters contains exactly two $1$s, so
$\nu_{3t+r}$ (the number of $1$s before position $3t + r$) equals
$2t$ plus the count inside the current block. By \eqref{eq:kernel},
if $w_t = 1$ the block is $(1,1,0)$, with $1$s at $r = 0$
($\nu = 2t$) and $r = 1$ ($\nu = 2t + 1$); if $w_t = 0$ the block is
$(0,1,1)$, with $1$s at $r = 1$ ($\nu = 2t$) and $r = 2$
($\nu = 2t + 1$). The block contributions to
$\sum_{w_i = 1} 2^i 3^{-(1+\nu_i)}$ are therefore
\[
w_t = 1:\;\; 2^{3t} 3^{-(2t+1)} + 2^{3t+1} 3^{-(2t+2)}
  = \Bigl(\tfrac{8}{9}\Bigr)^{t} \cdot \tfrac{5}{9}, \qquad
w_t = 0:\;\; 2^{3t+1} 3^{-(2t+1)} + 2^{3t+2} 3^{-(2t+2)}
  = \Bigl(\tfrac{8}{9}\Bigr)^{t} \cdot \tfrac{10}{9}.
\]
Grouping consecutive terms of a convergent series in a complete
ultrametric field is harmless, so
\[
x(w) = -\sum_{t \ge 0} \Bigl(\tfrac{8}{9}\Bigr)^{t}
  \Bigl(\tfrac{10}{9} - \tfrac{5}{9} w_t\Bigr)
  = -\tfrac{10}{9} \cdot \frac{1}{1 - \tfrac89}
    + \tfrac{5}{9} \sum_{t\ge0} w_t \Bigl(\tfrac{8}{9}\Bigr)^{t}
  = -10 + \tfrac59\, F_2(8/9). \qedhere
\]
\end{proof}

\begin{remark}
Iterating Proposition~\ref{prop:feq} at $z = 8/9$ unrolls
Theorem~\ref{thm:eval} into a tower: with
$u_k := \beta_k^{-1} F_2(\alpha_k/\beta_k)$,
$(\alpha_k, \beta_k) = (8^{3^k}\!, 9^{3^k})$,
\[
u_k \;=\; \frac{\alpha_k(\alpha_k + \beta_k)}{\beta_k^3 - \alpha_k^3}
\;+\; (\beta_k^2 - \alpha_k^2)\, u_{k+1},
\qquad\text{e.g.}\quad u_0 = \tfrac{136}{217} + 17\, u_1,
\]
with constants of height $\approx 9^{3^{k+1}}$. The tower view makes
the height growth of the method visible; the single-equation view is
the one the classical theory consumes.
\end{remark}

\section{The reduction}\label{sec:reduction}

\begin{corollary}[rung 1 for $w$, reduced]\label{cor:reduction}
If $F_2(8/9) \notin \Q$, then no positive integer has itinerary $w$.
\end{corollary}

\begin{proof}
If $n \in \Z_{\ge1}$ had itinerary $w$, Proposition
\ref{prop:realization} and Theorem~\ref{thm:eval} give
$F_2(8/9) = \tfrac95(n + 10) \in \Q$.
\end{proof}

The same reduction holds across the natural class.

\begin{proposition}\label{prop:class}
Let $\sigma$ be a substitution on $\{0,1\}$ of constant length $L$
whose two image blocks each contain exactly $a$ ones, with
$2^L < 3^a$ (equivalently: the fixed-point density $a/L$ exceeds
$\delta_c$), and let $w$ be a fixed point with generating function
$F$. Then $F(z) = A_\sigma(z) + B_\sigma(z) F(z^L)$ with
$A_\sigma, B_\sigma \in \Q(z)$ explicitly computable from the two
blocks, and there are rationals $A, B$ (computable from the blocks;
$B \neq 0$ unless the two blocks contribute identically, in which case
$x(w) \in \Q$ and the question is decidable directly) with
\[
x(w) \;=\; A + B \cdot F_2\!\left(\frac{2^L}{3^a}\right),
\]
the series converging both $2$-adically
($|2^L/3^a|_2 = 2^{-L}$) and in $\R$ ($0 < 2^L/3^a < 1$, by
supercriticality). Hence irrationality of the single $2$-adic Mahler
value $F_2(2^L/3^a)$ excludes all integers from realizing $w$.
\end{proposition}

\begin{proof}
Identical to Proposition~\ref{prop:feq} and Theorem~\ref{thm:eval}:
blockwise grouping. Each block of $L$ letters contains exactly $a$
ones, so $\nu_{Lt + r} = at + (\text{count inside the block})$, the
type-$b$ block ($b \in \{0,1\}$) contributes
$z^t C_b$ with $z = 2^L/3^a$ and
$C_b = \sum_{r:\, \sigma(b)_r = 1} 2^r 3^{-(1 + \#\{r' < r:\,
\sigma(b)_{r'} = 1\})}$, and
$x(w) = -\sum_t z^t (C_0 + (C_1 - C_0) w_t)
= -C_0/(1 - z) - (C_1 - C_0) F_2(z)$.
For the functional equation, group $F$ in $L$-blocks as before. The
double convergence is the displayed absolute-value computation;
$3^a > 2^L$ is exactly $z < 1$ in $\R$.
\end{proof}

The supercritical band --- the only band where divergent orbits can
live --- is thus exactly the band where the Mahler series converges at
\emph{both} places. We do not currently use the archimedean
convergence, but the classical method does (Section~\ref{sec:route}),
and the coincidence deserves to be on the record.

\section{The gap: truncation provably cannot decide}\label{sec:gap}

Why is this not finished? The series defining $F_2(8/9)$ converges
fast, and fast rational approximation is the standard engine of
irrationality proofs. The following proposition shows the engine
stalls --- by a precisely measurable margin.

\begin{proposition}\label{prop:gap}
For $N \ge 1$ let $P_N/Q_N$ with
$P_N = \sum_{i<N} w_i\, 8^i\, 9^{N-1-i} \in \Z$, $Q_N = 9^N$, be the
$N$-term truncation of $F_2(8/9)$. Then:
\begin{enumerate}
\item (quality) $\bigl|F_2(8/9) - P_N/Q_N\bigr|_2 = 2^{-3 i_0(N)}$
where $i_0(N) \in \{N, N+1, N+2\}$ is the first index $\ge N$ with
$w_{i_0} = 1$; in particular the error is never $0$ and
\[
\bigl|F_2(8/9) - P_N/Q_N\bigr|_2
\;=\; Q_N^{-\kappa_N}, \qquad
\kappa_N \;\to\; \frac{3\log 2}{2 \log 3} \;=\; 0.946394\ldots
\]
\item (insufficiency) For every rational $\xi = a/b$ and every
rational $P/Q \neq \xi$ with $Q$ odd and $|P| \le Q$,
\[
\bigl|\xi - P/Q\bigr|_2 \;\ge\; \frac{1}{(|a| + |b|)\, Q}.
\]
Consequently the approximations of part (1) are consistent with
$F_2(8/9) \in \Q$: any family of rational approximations can certify
irrationality only with exponent $> 1$, and the truncations achieve
$0.9464$.
\end{enumerate}
\end{proposition}

\begin{proof}
(1) The error is $\sum_{i \ge N} w_i (8/9)^i$; its terms have
strictly increasing $2$-adic valuations $3i$, so the sum is nonzero
with absolute value that of its first nonzero term, $2^{-3 i_0}$. The
word $w$ contains no three consecutive zeros (every block has a $1$ in
the middle position, by \eqref{eq:kernel}), so $i_0(N) \le N + 2$.
Then $\kappa_N = 3 i_0(N) \log 2 / (2N \log 3) \to (3\log2)/(2\log3)$.
(2) $\xi - P/Q = (aQ - bP)/(bQ)$; the numerator is a nonzero integer,
so $|aQ - bP|_2 \ge |aQ - bP|^{-1} \ge ((|a|+|b|)Q)^{-1}$, while
$|bQ|_2 \le 1$.
\end{proof}

So the question is genuinely beyond the convergence speed of its own
defining series: the gap between $0.9464$ and $1$ must be crossed by
structure, not by speed. Note where the constant comes from. For a
general admissible pair $(L, a)$ (Proposition~\ref{prop:class}), the
$N$-term truncation gains $LN \log 2$ bits of $2$-adic smallness
against a denominator of $aN \log 3$ bits, so the exponent is
\[
\kappa \;=\; \frac{L \log 2}{a \log 3}
\;=\; \frac{\delta_c}{\,a/L\,},
\]
the critical density divided by the word's density --- and the
supercriticality condition $a/L > \delta_c$ is \emph{exactly} the
statement $\kappa < 1$, with equality on the critical line. The
constant that creates the divergence problem also protects this
corner of it: throughout the band where divergent orbits could live,
truncation falls short of irrationality by construction.

\section{The route: Mahler's method, and one open question}
\label{sec:route}

Mahler's method \cite{mahler, nishioka} was invented for precisely
the situation of Section~\ref{sec:gap}: values of functions
satisfying $F(z) = A(z) + B(z) F(z^d)$ at algebraic points, where
naive approximation falls short. Its auxiliary-polynomial
construction converts the functional equation into approximation
families whose exponent grows with the auxiliary degree --- crossing
any fixed Liouville threshold --- and its hypotheses are, for our
instance, verified above: $F$ transcendental over $\Q(z)$
(Lemma~\ref{lem:fatou}), the point $8/9$ algebraic with orbit
$(8/9)^{3^k}$ avoiding the singularities of the coefficients (the
poles of $(z+z^2)/(1-z^3)$ at cube roots of unity and the zeros of
$1 - z^2$ at $\pm 1$).

At the archimedean place the classical theorem applies as stated:

\begin{theorem}\label{thm:arch}
The real number $F(8/9) = \sum_i w_i (8/9)^i \in \R$ is
transcendental.
\end{theorem}

\begin{proof}
Mahler's theorem (\cite{mahler}; in the form
\cite[Ch.~1]{nishioka}): if $F \in \Q[[z]]$ converges in
$|z| < 1$, satisfies a functional equation of the class above, and is
transcendental over $\Q(z)$, then $F(\alpha)$ is transcendental for
every algebraic $\alpha$, $0 < |\alpha| < 1$, whose orbit
$\alpha^{d^k}$ avoids the finitely many excluded points (poles and
zeros of the coefficient functions). Lemma~\ref{lem:fatou} and the
singularity check above verify every hypothesis at $\alpha = 8/9$.
\end{proof}

The number rung 1 needs is the \emph{same} function at the
\emph{same} rational point, evaluated at the other place where the
series converges:

\begin{question}\label{q:main}
Is $F_2(8/9) \in \Qtwo$ irrational? More generally: does Mahler's
method, transposed to $\Qtwo$ --- auxiliary polynomials, the
non-archimedean Schwarz lemma in place of the maximum principle,
heights as before --- yield irrationality (or transcendence) of
$F_p(\alpha)$ for Mahler functions $F \in \Q[[z]]$ at rational
$\alpha$ with $0 < |\alpha|_p < 1$ under the hypotheses of
Theorem~\ref{thm:arch}?
\end{question}

We are not aware of a published theorem in exactly this form, though
$p$-adic and function-field variants of Mahler's method exist in
several settings (see \cite{nishioka} and its references;
\cite{fernandes} develops the method over function fields in positive
characteristic, and \cite{af} raises the structural theory at the
archimedean place to a complete description of algebraic relations).
Locating the statement, or transposing the archimedean proof --- the
ingredients are place-generic, and non-archimedean analysis is, if
anything, kinder to the relevant lemmas --- is the concrete work item
this paper isolates. We record what is at stake: a positive answer to
Question~\ref{q:main} excludes, by Corollary~\ref{cor:reduction} and
Proposition~\ref{prop:class} (plus a finite-system extension to
shifted fixed points, standard in the Mahler literature \cite{af}),
the entire uniform-substitution stratum of the automatic rung.

Two further remarks on the interface. First, the celebrated
transcendence results for automatic \emph{numbers}
\cite{ab} do not apply here: $x(w)$ is not an automatic number ---
its binary digits are not generated by an automaton; only the
\emph{series coefficients} are automatic, and the base mixes powers
of $2$ against powers of $3$. The functional equation is the correct
interface, not the digit expansion. Second, the obstruction is not
the word combinatorics: all of Sections
\ref{sec:realization}--\ref{sec:gap} are exact and verified; the
single missing ingredient is non-archimedean.

\section{Calibration: the $5x+1$ test}\label{sec:calibration}

Any principle that would exclude divergence must fail to ``prove''
the same for the $5x+1$ map, where orbits (conjecturally, and
overwhelmingly in experiment, e.g.\ from $n = 7$) really do diverge.
The present method passes this filter in the correct way. For $5x+1$,
the identical computation realizes the same word $w$ at
$x_5(w) = A' + B' F_2(2^L/5^a)$ --- the same generating function at a
different rational point ($8/25$ for our $\sigma$) --- and a $2$-adic
Mahler argument would conclude the same way: \emph{no integer has
that exact automatic itinerary under $5x+1$}. This is consistent with
(indeed expected alongside) real divergence: the observed $5x+1$
escapees have generic-looking itineraries of density near $0.49$,
nowhere near a substitution fixed point. The method excludes specific
symbolic classes; it cannot manufacture global convergence --- which
is precisely why it can be sound. The $3x+1$-specific content remains
where the companion papers put it: in the density floor that pins
divergent orbits into the narrow band where the ordered words live.

\section{Scope, and the ladder above}\label{sec:scope}

What would closing rung 1 actually buy? Directly: no divergent orbit
has an itinerary that is a (shifted) fixed point of a uniform
substitution --- a countable family, measure zero in every sense.
The conjecture would remain open at every higher rung: morphic words,
Sturmian words (where the companion shadow paper locates the
slowest-escape candidates), and the unstructured majority, which the
no-go theorem of \cite{companion} places beyond all bounded-memory
arguments. The honest value of the rung is the bridge it builds: an
exact, mostly machine-verified chain from the Collatz dynamics to a
single value of a classical Mahler function, with the remaining gap
isolated as one well-posed question (Question~\ref{q:main}) in a
mature theory --- and a quantified certificate
(Proposition~\ref{prop:gap}) that the question sits genuinely beyond
elementary methods, for the problem's own critical reason.

\subsection*{Reproducibility}\label{sec:repro}

The formal results (\code{eq\_of\_itinerary\_eq},
\code{divergent\_itinerary\_aperiodic}, \code{dcoef\_closed},
\code{terras\_exact\_form}, the parity bijection) are in
\code{lean/Collatz/Ladder.lean}, \code{Shadow.lean}, and
\code{Parity.lean} of \cite{companion} (Lean~4 + mathlib; no axioms
beyond the standard three, no \code{sorry}). The numerical
verifications --- the functional equation to degree $900$, the
evaluation identity to $600$ bits of $2$-adic precision, the exact
finite-level tower identities, the block law on $200$ random words,
the realization digit-tail statistics, and the shadow clustering of
automatic words --- are
\code{analysis/mahler\_experiments.py} (experiments 1--4),
deterministic, run time under a minute. The attack plan in extended
form, including the failed routes, is
\code{analysis/RUNG1\_ATTACK.md}.

\subsection*{Acknowledgments}

The mathematics, the formalization, and this text were produced in
collaboration with Claude (Anthropic). The reduction owes its shape
to the exact-form and cocycle infrastructure of the companion
formalization; the question it isolates belongs to a theory built by
Mahler and rebuilt by many hands since.

\begin{thebibliography}{10}

\bibitem{ab} B.~Adamczewski and Y.~Bugeaud, On the complexity of
algebraic numbers I. Expansions in integer bases, \emph{Ann. of
Math.} \textbf{165} (2007), 547--565.

\bibitem{af} B.~Adamczewski and C.~Faverjon, M\'ethode de Mahler:
relations lin\'eaires, transcendance et applications aux nombres
automatiques, \emph{Proc. Lond. Math. Soc.} \textbf{115} (2017),
55--90.

\bibitem{as} J.-P.~Allouche and J.~Shallit, \emph{Automatic
Sequences: Theory, Applications, Generalizations}, Cambridge
University Press, 2003.

\bibitem{fatou} P.~Fatou, S\'eries trigonom\'etriques et s\'eries de
Taylor, \emph{Acta Math.} \textbf{30} (1906), 335--400.

\bibitem{fernandes} G.~Fernandes, M\'ethode de Mahler en
caract\'eristique non nulle, \emph{Ann. Inst. Fourier} \textbf{68}
(2018).

\bibitem{lagarias-overview} J.~C.~Lagarias, The $3x+1$ problem: an
overview, arXiv:2111.02635 (2021).

\bibitem{lvdp} J.~H.~Loxton and A.~J.~van der Poorten, Arithmetic
properties of the solutions of a class of functional equations,
\emph{J. Reine Angew. Math.} \textbf{330} (1982), 159--172.

\bibitem{mahler} K.~Mahler, Arithmetische Eigenschaften der
L\"osungen einer Klasse von Funktionalgleichungen, \emph{Math. Ann.}
\textbf{101} (1929), 342--366.

\bibitem{nishioka} Ku.~Nishioka, \emph{Mahler Functions and
Transcendence}, Lecture Notes in Mathematics 1631, Springer, 1996.

\bibitem{companion} M.~Sharpe, The Collatz conjecture at the critical
line: machine-verified no-go, density, and cycle theorems, companion
papers and repository, 2026.

\bibitem{terras} R.~Terras, A stopping time problem on the positive
integers, \emph{Acta Arith.} \textbf{30} (1976), 241--252.

\end{thebibliography}

\end{document}
